% ===========================================================================
\chapter*{Conclusion~: le voyage complet}
\addcontentsline{toc}{chapter}{Conclusion}
% ===========================================================================

\begin{maitre}
R\'esumons notre voyage. Qu'est-ce que fait ce logiciel,
dans tes propres mots~?
\end{maitre}

\begin{apprenti}
Il lit un stock de poutres bois, g\'en\`ere toutes les configurations
de calcul (en faisant le produit cart\'esien des param\`etres),
et pour chaque poutre et chaque port\'ee, il v\'erifie la r\'esistance
(ELU~: flexion, cisaillement, appui, d\'eversement) et les d\'eformations
(ELS~: fl\`eches instantan\'ee et finale) selon l'Eurocode~5.
Il stocke les r\'esultats dans un CSV et m\'emorise ce qui a d\'ej\`a
\'et\'e calcul\'e pour ne pas tout recalculer \`a chaque fois.
\end{apprenti}

\begin{maitre}
Et Python dans tout \c{c}a~?
\end{maitre}

\begin{apprenti}
Python orchestre tout. Les \textbf{variables} stockent les donn\'ees
(avec l'unit\'e dans le nom). Les \textbf{fonctions} encapsulent les formules
(une fonction = une formule de la norme). Les \textbf{classes} mod\'elisent
les types de poutres avec leur comportement propre~: c'est le polymorphisme.
Les \textbf{boucles} parcourent le stock et les port\'ees. NumPy
\textbf{vectorise} les calculs pour la performance. Pandas lit et \'ecrit
les CSV. Et Pydantic valide les donn\'ees d'entr\'ee.
\end{apprenti}

\begin{important}
Les concepts Python rencontr\'es dans ce logiciel~--- variables, fonctions,
classes, polymorphisme, boucles, dictionnaires, modules~--- sont exactement
ceux qu'on enseigne en premi\`ere ann\'ee de programmation.
La diff\'erence, c'est qu'ici ils sont mis au service d'un probl\`eme
d'ing\'enierie r\'eel et norm\'e.
La structure du code \textbf{refl\`ete} la structure de la norme~:
un module \py{ec0/} pour EN~1990, un module \py{ec1/} pour EN~1991,
et un module \py{ec5/} pour EN~1995.
Ce n'est pas un hasard~: c'est une conception d\'elib\'er\'ee.
\end{important}

\vspace{2em}
\begin{center}
\rule{0.5\textwidth}{0.5pt}
\end{center}
\vspace{1em}

\begin{tcolorbox}[
  colback=bleuMaitreLight, colframe=bleuMaitre,
  width=0.9\textwidth, center
]
\centering\itshape
\guillemotleft\, La connaissance s'acqui\`ere par l'exp\'erience,\\
tout le reste n'est qu'information.\,\guillemotright\\[4pt]
{\normalsize --- Albert Einstein}
\end{tcolorbox}